\documentclass[11pt]{article} 
\usepackage[margin=1in]{geometry}
\usepackage{amsfonts,latexsym,amsthm,amssymb,amsmath,amscd,euscript,esint, dsfont,mathtools,stmaryrd}	
\usepackage{graphicx}		
\usepackage{hyperref}
\usepackage{cases}			
\usepackage{textcomp, gensymb}
\usepackage{xcolor}
\usepackage{tabularx}

% diagrams
\usepackage{quiver}

\makeatletter
\def\namedlabel#1#2{\begingroup
	\def\@currentlabel{#2}
	\label{#1}\endgroup
}
\makeatother

\setlength{\parindent}{0pt} 	% first line in paragraph will not be indented

\usepackage[parfill]{parskip}
\usepackage{lastpage}
\usepackage{fancyhdr}
\newcommand{\ind}{{\mathds{1}}}

% math fonts
\renewcommand{\AA}{\mathbb{A}}
\newcommand{\BB}{\mathbb{B}}
\newcommand{\CC}{\mathbb{C}}
\newcommand{\DD}{\mathbb{D}}
\newcommand{\EE}{\mathbb{E}}
\newcommand{\FF}{\mathbb{F}}
\newcommand{\GG}{\mathbb{G}}
\newcommand{\HH}{\mathbb{H}}
\newcommand{\II}{\mathbb{I}}
\newcommand{\JJ}{\mathbb{J}}
\newcommand{\KK}{\mathbb{K}}
\newcommand{\LL}{\mathbb{L}}
\newcommand{\MM}{\mathbb{M}}
\newcommand{\NN}{\mathbb{N}}
\newcommand{\OO}{\mathbb{O}}
\newcommand{\PP}{\mathbb{P}}
\newcommand{\QQ}{\mathbb{Q}}
\newcommand{\RR}{\mathbb{R}}
\renewcommand{\SS}{\mathbb{S}}
\newcommand{\TT}{\mathbb{T}}
\newcommand{\UU}{\mathbb{U}}
\newcommand{\VV}{\mathbb{V}}
\newcommand{\WW}{\mathbb{W}}
\newcommand{\XX}{\mathbb{X}}
\newcommand{\YY}{\mathbb{Y}}
\newcommand{\ZZ}{\mathbb{Z}}

\newcommand{\cA}{\mathcal{A}}
\newcommand{\cB}{\mathcal{B}}
\newcommand{\cC}{\mathcal{C}}
\newcommand{\cD}{\mathcal{D}}
\newcommand{\cE}{\mathcal{E}}
\newcommand{\cG}{\mathcal{G}}
\newcommand{\cF}{\mathcal{F}}
\newcommand{\cH}{\mathcal{H}}
\newcommand{\cI}{\mathcal{I}}
\newcommand{\cJ}{\mathcal{J}}
\newcommand{\cK}{\mathcal{K}}
\newcommand{\cL}{\mathcal{L}}
\newcommand{\cM}{\mathcal{M}}
\newcommand{\cN}{\mathcal{N}}
\newcommand{\cO}{\mathcal{O}}
\newcommand{\cP}{\mathcal{P}}
\newcommand{\cQ}{\mathcal{Q}}
\newcommand{\cR}{\mathcal{R}}
\newcommand{\cS}{\mathcal{S}}
\newcommand{\cT}{\mathcal{T}}
\newcommand{\cU}{\mathcal{U}}
\newcommand{\cV}{\mathcal{V}}
\newcommand{\cW}{\mathcal{W}}
\newcommand{\cX}{\mathcal{X}}
\newcommand{\cY}{\mathcal{Y}}
\newcommand{\cZ}{\mathcal{Z}}

\newcommand{\ka}{\mathfrak{a}}
\newcommand{\kb}{\mathfrak{b}}
\newcommand{\kc}{\mathfrak{c}}
\newcommand{\kd}{\mathfrak{d}}
\newcommand{\ke}{\mathfrak{e}}
\newcommand{\kg}{\mathfrak{g}}
\newcommand{\kf}{\mathfrak{f}}
\newcommand{\kh}{\mathfrak{h}}
\newcommand{\ki}{\mathfrak{i}}
\newcommand{\kj}{\mathfrak{j}}
\newcommand{\kk}{\mathfrak{k}}
\newcommand{\kl}{\mathfrak{l}}
\newcommand{\km}{\mathfrak{m}}
\newcommand{\kn}{\mathfrak{n}}
\newcommand{\ko}{\mathfrak{o}}
\newcommand{\kp}{\mathfrak{p}}
\newcommand{\kq}{\mathfrak{q}}
\newcommand{\kr}{\mathfrak{r}}
\newcommand{\ks}{\mathfrak{s}}
\newcommand{\kt}{\mathfrak{t}}
\newcommand{\ku}{\mathfrak{u}}
\newcommand{\kv}{\mathfrak{v}}
\newcommand{\kw}{\mathfrak{w}}
\newcommand{\kx}{\mathfrak{x}}
\newcommand{\ky}{\mathfrak{y}}
\newcommand{\kz}{\mathfrak{z}}

\newcommand{\kA}{\mathfrak{A}}
\newcommand{\kB}{\mathfrak{B}}
\newcommand{\kC}{\mathfrak{C}}
\newcommand{\kD}{\mathfrak{D}}
\newcommand{\kE}{\mathfrak{E}}
\newcommand{\kG}{\mathfrak{G}}
\newcommand{\kF}{\mathfrak{F}}
\newcommand{\kH}{\mathfrak{H}}
\newcommand{\kI}{\mathfrak{I}}
\newcommand{\kJ}{\mathfrak{J}}
\newcommand{\kK}{\mathfrak{K}}
\newcommand{\kL}{\mathfrak{L}}
\newcommand{\kM}{\mathfrak{M}}
\newcommand{\kN}{\mathfrak{N}}
\newcommand{\kO}{\mathfrak{O}}
\newcommand{\kP}{\mathfrak{P}}
\newcommand{\kQ}{\mathfrak{Q}}
\newcommand{\kR}{\mathfrak{R}}
\newcommand{\kS}{\mathfrak{S}}
\newcommand{\kT}{\mathfrak{T}}
\newcommand{\kU}{\mathfrak{U}}
\newcommand{\kV}{\mathfrak{V}}
\newcommand{\kW}{\mathfrak{W}}
\newcommand{\kX}{\mathfrak{X}}
\newcommand{\kY}{\mathfrak{Y}}
\newcommand{\kZ}{\mathfrak{Z}}

%some math symbol commands redefined:
\DeclareMathOperator{\C}{\mathcal{C}}
\DeclareMathOperator{\power}{\mathcal{P}}
\DeclareMathOperator{\vspan}{\text{span}}
\DeclareMathOperator{\vnull}{\text{null}}
\DeclareMathOperator{\range}{\text{range}}
\DeclareMathOperator{\re}{\text{Re}}
\DeclareMathOperator{\im}{\text{Im}}
\DeclareMathOperator{\erf}{\text{erf}}
\newcommand{\pdv}[2]{\frac{\partial #1}{\partial #2}}
\newcommand{\pdvv}[2]{\frac{\partial^2 #1}{\partial #2}}
\DeclareMathOperator{\tr}{\operatorname{Tr}}
\newcommand{\Deck}{\operatorname{Deck}}
\newcommand{\Conj}{\mathrm{Conj}}
\newcommand{\Sing}{\mathrm{Sing}}
\DeclareMathOperator{\Spec}{\operatorname{Spec}}
\newcommand{\Proj}{\operatorname{Proj}}
\newcommand{\Res}{\operatorname{Res}}
\newcommand{\PROJ}{\mathit{Proj}}

% category names
\newcommand{\Set}{\mathsf{Set}}
\newcommand{\Grp}{\mathsf{Grp}}
\newcommand{\Ring}{\mathsf{Ring}}
\newcommand{\Top}{\mathsf{Top}}
\newcommand{\Sch}{\mathsf{Sch}}

\newcommand{\ob}{\operatorname{Ob}}
\newcommand{\Id}{\operatorname{Id}}
\newcommand{\Hom}{\operatorname{Hom}}
\newcommand{\colim}{\operatorname{colim}}
\newcommand{\Ann}{\operatorname{Ann}}
\newcommand{\Supp}{\operatorname{Supp}}
\newcommand{\Ass}{\operatorname{Ass}}
\newcommand{\pre}{\text{pre}}
\newcommand{\adj}{\text{adj}}
\newcommand{\Ker}{\operatorname{Ker}}
\newcommand{\Coker}{\operatorname{Coker}}
\newcommand{\Nil}{\operatorname{Nil}}
\newcommand{\Char}{\operatorname{char}}
\newcommand{\Adj}{\operatorname{Adj}}
\newcommand{\Bl}{\mathrm{Bl}}
\newcommand{\codim}{\mathrm{codim}}
\newcommand{\val}{\operatorname{val}}
\newcommand{\Div}{\operatorname{div}}
\newcommand{\Pic}{\operatorname{Pic}}
\newcommand{\Weil}{\operatorname{Weil}}
\newcommand{\Cl}{\operatorname{Cl}}
\newcommand{\Sym}{\operatorname{Sym}}
\newcommand{\Aut}{\operatorname{Aut}}
\newcommand{\Gr}{\operatorname{Gr}}

\newcommand{\GL}{\operatorname{GL}}
\newcommand{\PGL}{\operatorname{PGL}}
\newcommand{\SL}{\operatorname{SL}}
\newcommand{\PSL}{\operatorname{PSL}}
\newcommand{\Rk}{\operatorname{Rk}}
\newcommand{\Tor}{\operatorname{Tor}}

\newtheorem{lemma}{Lemma}
\newtheorem{claim}{Claim}

% category theory symbols
\DeclareMathOperator{\Mor}{\operatorname{Mor}}

\newenvironment{solution}{\begin{trivlist}\item[]{\textcolor{blue}{Solution:}}}{\qed \end{trivlist}}


\usepackage[shortlabels]{enumitem}
\title{MATH 2060 Homework 12}
\author{Zhengyu Zou}
\date{\today}

\begin{document}
\maketitle

\textbf{Collaborators:} Ethan Bove, Roberta Pagliaro, Anamaria Perez

\section*{FOAG 24.3.D.}

\begin{solution}
    Denote by $A = k[t]/(t^2)$. 
    We only need to check that $\Tor_1^{A} (M, A / I) = 0$ for all ideals $I \subset A$. Notice that $A$ has only three ideals: $(0)$, $(t)$, and $(1)$. Let $M_i$ be a projective resolution of $M$. Clearly, $A / (1) = 0$, and $Tor_1^A (M, 0) = 0$. Also, since $M_2 \rightarrow M_1 \rightarrow M_0 \rightarrow M \rightarrow 0$ is exact, and $M_i \otimes_A A = M_i$, we see that $\Tor_1^A(M, A / (0)) = \Tor_1^A(M, A) = 0$. Finally, recall that $\Tor_1^A (M, A/(t)) = \Tor_1^A (A/(t), M)$. Consider a free resolution of $A/(t)$ given by 
    \[
        \cdots \xrightarrow{\times t} \frac{k[t]}{(t^2)} \xrightarrow{\times t} \frac{k[t]}{(t^2)} \xrightarrow{\times t} \frac{k[t]}{(t^2)} \rightarrow k \rightarrow 0.
    \]
    We then tensor with $M$ to get 
    \[
        \cdots \xrightarrow{\times t} M \xrightarrow{\times t} M \xrightarrow{\times t} M \rightarrow 0.
    \]
    We see that $\Tor_1^A(M, A/(t))$ vanishes iff the kernel of the map $M \xrightarrow{\times t} M$ equals its image, which is true iff $M / tM \xrightarrow{\times t} tM$ is an isomorphism. This completes the proof. 
\end{solution}

\newpage

\section*{FOAG 24.3.H.}

\begin{solution}
    Let $X = \Spec k[t]$ and $Y = \Spec k[x,y]/(y^2 - x^2 - x^3)$.
    The normalization of the node is given by the map $X \rightarrow Y$ sending $(x, y) \mapsto (t^2 - 1, t^3 - t)$. The preimage of any point $q = [(x-a,y-b)]$ with $(a,b) \neq (0,0)$ is a single point, so the dimension of the stalk $(\pi_* \cO_X)_q \times \kappa(q)$ is 1 (notice that $X_{t^2 - 1}$ is isomorphic to $Y_{x}$ by sending $t \mapsto \frac{y}{x}$). On the other hand, the preimage of $q = [(x,y)]$ consists of two points $[(t-1)], [(t+1)]$, and $\cO_X$ is nontrivial on both points. This implies $(\pi_* \cO_X)_q \times \kappa(q) \geq 2$, but then that means the rank of $\pi_* \cO_X$ is not locally constant. 
\end{solution}

\newpage

\section*{FOAG 24.3.L.}

\begin{solution}
\begin{enumerate}[(a)]
    \item We know that the family of pairs of points is the vanishing of the ideal $(x - t, y - 4t) \cdot (x, y)$ in $\AA^3$, which is cut out by four equations $x(x - 3t)$, $x (y - 4t)$, $y (x - 3t)$, and $y (y - 4t)$. Away from $t = 0$, we see that $t$ is invertible, so the scheme-theoretical closure of the complement of $t = 0$ is cut out by the extra equation $3xt - 4yt$, but since we are allowed to divide by $t$, we have the equation $3x - 4y$. The fiber is then cut out by four equations:
    \[
        x^2; \hspace{3pt} y^2; \hspace{3pt} xy; \hspace{3pt} 3x - 4y.
    \]
    Notice that we get $y^2$ and $xy $ from $3x - 4y$, so the flat limit is cut out by $x^2$ and $3x - 4y$. 

    \item To find the flat limit as $t \rightarrow \infty$, we would need to work in $\PP^2$ and parametrize the points with $\PP^1$. Consider the pair of points $\{[0:0:1], [3t:4t:1]\}$ in $\PP^2$ parametrized by $\PP^1$, we may switch parameters $s = t^{-1}$ to get a parametrization $\{[0:0:1], [3:4:s]\}$, and we would like to find the flat limit as $s \rightarrow 0$. It is then obvious that the flat limit is cut out by the product of the homogeneous ideal $(x, y)$ with $(z, 3x - 4y)$ in $\PP^2$, which is the disjoint union of two points $[0:0:1]$ and $[3:4:0]$. 
\end{enumerate}
\end{solution}

\newpage

\section*{FOAG 24.6.A.}

\begin{solution}
    Denote by $X \subset \AA^4$ the subscheme cut out by $(w, x) \cdot (y, z)$, and let $B$ denote its structure ring. 
    Consider the nonzero-divisor $a \in k[a,b]$, which gets mapped to $t = w - y$ in $X$. Notice that $V(t)$ does not contain the plane $V(w, x)$ or $V(y,z)$, so it must intersect each plane at a 1-dimensional subspace. We claim that $V(t)$ is not flat over $\Spec k[b]$. To see this, since $\Spec k[b]$ is a regular Noetherian curve, it suffices to show that there exists an associated point of $V(t)$ that maps to some closed point of $\Spec k[b]$. Consider the preimage of $(0, 0) \in \AA^2$, which is the point $(0, 0, 0, 0) \in \AA^4$. The Zariski tangent space of $(0,0, 0, 0)$ is 4-dimensional. Then, suppose $t$ is a nonzero-divisor, then it's an effective Cartier divisor, and the Zariski tangent space of $V(t)$ at $(0, 0, 0, 0)$ must have dimension 3. Now, if $V(t)$ contains no embedded point at $(0, 0, 0, 0)$, then it must be the scheme-theoretic intersection of the two 1-dimensional subspaces $V(t) \cap V(w,x)$ and $V(t) \cap V(y,z)$, but then it would have a 2-dimensional Zariski tangent space. Therefore, $V(t)$ must have an embedded point at $(0, 0, 0, 0)$, but then the map $V(t) \rightarrow \Spec k[b]$ is not flat, contradicting the slicing criterion. 
\end{solution}

\newpage

\section*{FOAG 24.7.F.}

\begin{solution}
    By symmetry, it suffices to show that $X$ is smooth over $U_1$. We know that if $X$ and $Y$ are both smooth and $X \hookrightarrow Y$ is a closed embedding, then $\Bl_X Y$ is smooth. Clearly, $C_1 \setminus \{p_1\}$ and $U_1$ are both smooth, and the map $C_1 \setminus \{p_1\} \hookrightarrow U_1$ is a closed embedding. This implies $\Bl_{C_1 \setminus \{p_1\}} U_1$ is smooth. Then, since the proper transform $C_2 \setminus \{p_1\} \hookrightarrow \Bl_{C_1 \setminus \{p_1\}} U_1$ is a closed embedding, we again see that $X_i$ is smooth. We conclude that $X$ is smooth. 

    Since $C_i \setminus \{p_1, p_2\} \hookrightarrow \PP^3$ is a regular embedding, 22.3.D(a) implies the exceptional divisor of $C_i \setminus \{p_1, p_2\}$ is the projectivized normal bundle. It is a $\PP^1$-bundle since $C_i \hookrightarrow \PP^3$ is smooth, and hence the normal bundle $N_{C_i} \PP^3$ is a locally free sheaf of rank 2. It is flat because the projectivized normal bundle $N_{C_i} \PP^3$ is smooth, which implies it's reduced and locally irreducible. It follows that there exists only one associated point on each component of $N_{C_i \setminus \{p_1, p_2\}} \PP^3$ (namely the generic point), which must map to the generic point of $C_i \setminus \{p_1, p_2\}$. By 24.3.J., the map $\pi|_{E_i}$ is flat. 
\end{solution}

\newpage

\section*{FOAG 24.7.G.}

\begin{solution}
    Since we have already shown $E_i$ is flat over $C_i \setminus \{p_1, p_2\}$, we only need to check flatness over $p_1$ and $p_2$. We can do so locally. Observe that on $U_i$, the scheme-theoretic closure of $E_i$ is simply the flat limit of $C_i$ at $p_i$, so it must be flat over $p_i$. This shows $\bar E_i \rightarrow C_i$ is flat. 
\end{solution}

\newpage

\section*{FOAG 24.7.H.}

\begin{solution}
    To see that $\pi^{-1}(p_{3-i})$ is two copies of $\PP^1$ glued along a point, we can work locally on $U_{i}$. Denote by $C'_j = C_j \setminus \{p_{i}\}$. Consider the map $\varphi: \Bl_{C'_{i}} U_{i} \rightarrow U_{i}$ obtained from blowing up $U_{i}$ along $C'_{i}$. Consider also $\psi: X_i \rightarrow \Bl_{C'_i} U_i$ obtained from blowing up $\Bl_{C'_i} U_i$ along the proper transform of $C'_{3-i}$. We see that $\varphi^{-1}(p_{3-i})$ is a single copy of $\PP^1$ that sits in the projective normal bundle of $C'_i$. On the other hand, blowing up along the proper transform of $C'_{3-i}$ picks up another copy of $\PP^1$, since the proper transform is a closed embedding into $\Bl_{C'_i} U_i$ and has codimension-2. The two copies of $\PP^1$ are glued together along the tangent direction of $C'_{3-i}$ at $p_{3-i}$. Denote by $Y_1$ the copy of $\PP^1$ at $p_1$ picked up from $C_1$, $Z_1$ the copy of $\PP^1$ at $p_1$ picked up from $C_2$. Define $Y_2, Z_2$ over $p_2$ similarly.

    To see that $\bar E_1$ contains only $Y_1, Y_2, Z_2$, notice that at $p_2$ we blow up $C_2$ and then $C_1$, so the flat limit of $C_1$ must pick up both copies of $\PP^1$, which corresponds to $Z_2$ and $Y_2$. On the other hand, at $p_1$ we blow up $C_1$ and then $C_2$, so the flat limit of $C_1$ does not pick up the copy of $\PP^1$ from $C_2$. Therefore, $Z_1$ is not in $\bar E_1$. 
\end{solution}

\newpage

\section*{FOAG 24.7.I.}

\begin{solution}
    Assume for contradiction that $X$ is projective. Let $\cL$ be a very ample line bundle over $X$. This implies $\cL$ has positive degree on $C_1, C_2, Y_1, Y_2, Z_1, Z_2$. Then, 24.7.D.\ implies 
    \[
        \deg_{Y_1} \cL = \deg_{Y_2} \cL + \deg_{Z_2} \cL ; \hspace{10pt}
        \deg_{Z_2} \cL = \deg_{Z_1} \cL + \deg_{Y_1} \cL.
    \]
    It follows that $\deg_{Y_2} \cL + \deg_{Z_1} \cL = 0$, but that contradicts the fact that both degrees are positive. 
\end{solution}

\newpage

\section*{Problem 9.}

Let $X = \Spec k[x, y]/(y^2 - x^3)$. Show that $X$ is not formally smooth over $\Spec k$
(directly from the definition not using the equivalence of smooth and formally smooth).

\begin{solution}
    Consider the following diagram 
    % https://q.uiver.app/#q=WzAsOCxbMCwwLCJcXFNwZWMgXFxmcmFje2tbdF19eyh0XjIpfSJdLFswLDEsIlxcU3BlYyBcXGZyYWN7a1t0XX17KHReMyl9Il0sWzEsMSwiXFxTcGVjIGsiXSxbMSwwLCJcXFNwZWMgXFxmcmFje2tbeCx5XX17KHleMiAtIHheMyl9Il0sWzIsMCwiXFxmcmFje2tbdF19eyh0XjIpfSJdLFszLDEsImsiXSxbMiwxLCJcXGZyYWN7a1t0XX17KHReMyl9Il0sWzMsMCwiXFxmcmFje2tbeCx5XX17KHleMiAtIHheMyl9Il0sWzAsMV0sWzEsMl0sWzMsMl0sWzAsM10sWzUsNl0sWzYsNF0sWzUsN10sWzcsNF1d
    \[\begin{tikzcd}
        {\Spec \frac{k[t]}{(t^2)}} & {X} & {\frac{k[t]}{(t^2)}} & {\frac{k[x,y]}{(y^2 - x^3)}} \\
        {\Spec \frac{k[t]}{(t^3)}} & {\Spec k} & {\frac{k[t]}{(t^3)}} & k
        \arrow[from=1-1, to=1-2]
        \arrow[from=1-1, to=2-1]
        \arrow[from=1-2, to=2-2]
        \arrow[from=1-4, to=1-3]
        \arrow[from=2-1, to=2-2]
        \arrow[from=2-3, to=1-3]
        \arrow[from=2-4, to=1-4]
        \arrow[from=2-4, to=2-3]
    \end{tikzcd}\]
    Here, the map $k[x,y]/(y^2 - x^3) \rightarrow k[t]/(t^2)$ is given by sending $(x,y) \mapsto (0, t)$. We will show that there exists no lift $\Spec k[t]/(t^3) \rightarrow X$ that makes the left diagram commute. Suppose we do have a lift that corresponds to a map of rings $k[x,y]/(y^2 - x^3) \rightarrow k[t]/(t^3)$. Let $a,b,c,d,e,f$ be elements of $k$ such that $x \mapsto a + bt + ct^2$ and $y \mapsto d + et + ft^2$. Since the following triangle commutes:
    % https://q.uiver.app/#q=WzAsMyxbMCwwLCJcXFNwZWMgXFxmcmFje2tbdF19eyh0XjIpfSJdLFswLDEsIlxcU3BlYyBcXGZyYWN7a1t0XX17KHReMyl9Il0sWzEsMCwiXFxTcGVjIFxcZnJhY3trW3gseV19eyh5XjIgLSB4XjMpfSJdLFswLDFdLFswLDJdLFsxLDJdXQ==
    \[\begin{tikzcd}
        {\Spec \frac{k[t]}{(t^2)}} & {\Spec \frac{k[x,y]}{(y^2 - x^3)}} \\
        {\Spec \frac{k[t]}{(t^3)}}
        \arrow[from=1-1, to=1-2]
        \arrow[from=1-1, to=2-1]
        \arrow[from=2-1, to=1-2]
    \end{tikzcd}\]
    we must have $a = b = d = 0$ and $e = 1$. It follows that $x \mapsto ct^2$ and $y \mapsto t + ft^2$. Then, since $y^2 - x^3 = 0$, we must have $(t + ft^2)^2 - (ct^2)^3 = 0$, but then $(t + ft^2)^2 = t^2 \neq 0$. This shows there exists no lift $\Spec k[t]/(t^3) \rightarrow X$. We conclude that $X$ is not formally smooth over $\Spec k$. 
\end{solution}

\newpage

\section*{FOAG 26.1.G.}

\begin{solution}
    We first observe that $R = k[w,x] \oplus k[y,z]$. It follows that elements of $R$ can be expressed uniquely as $f(w,x) + g(y,z) + c$, where $f(0,0) = g(0,0) = 0$, and $c \in k$. Now, suppose we have some $a = f(w,x) + g(y,z) + c \in R$ such that $(w - y) a = 0$. It follows that 
    \[
        (w - y) a 
        = w f(w , x) - y g (y, z) + (w - y)c = 0.
    \]
    Notice that the degrees of the terms in $wf(w,x)$ and $yg(y,z)$ are both > 1, so we must have $f(w,x) = g(y,z) \equiv 0$ and $c = 0$. It follows that $w - y$ is not a zero-divisor in $R$. In particular, $V(w-y)$ is an effective Cartier divisor of $X = \Spec R$. 

    The argument from 24.6.A.\ then shows that $\Spec R/(w-y)$ has an embedded point at $(0,0,0,0)$. To make it self-contained, we will reiterate the proof here: The projection map $X \rightarrow \AA^2$ given by the two coordinates $(a,b) \mapsto (w - y, x - z)$ is not flat. Therefore, by the slicing criterion, we know that the map $\Spec R / (w - y) \rightarrow \AA^1$ is not flat. Since it is flat at $\AA^1 \setminus \{0\}$, it must not be flat at $0$, but the only possible way for the map to not be flat at 0 is for $\Spec R/(w-y)$ to have an associated point at the origin $(0,0,0,0)$ that maps to $0 \in \AA^1$. This shows that the origin is an embedded point, so the depth of the local ring at the origin is 1. 
\end{solution}

\newpage

\section*{FOAG 26.2.E.}

\begin{solution}
    To show that $h \in (f,g)$, it suffices to show that $X = \Spec \CC[x_0, x_1, x_2] / (f, g)$ is reduced. That is, $X$ is Cohen-Macaulay at all closed points. Let $p$ be a closed point of $X$. Clearly, $\AA_{\CC}^3$ is Cohen-Macaulay at $p$, and since $\CC[x_0, x_1, x_2]$ is a UFD, $f$ is not a zero-divisor. It follows from the slicing criterion that $\Spec \CC[x_0, x_1, x_2] / (f)$ is Cohen-Macaulay at $p$. Next, to show that $X$ is Cohen-Macaulay at $p$, we only need to show that $g$ is not a zero-divisor in $\CC[x_0, x_1, x_2] / (g)$. This simply follows from the fact that $f$ and $g$ forms a complete intersection (the two curves intersect at a finite number of reduced points, so $(f,g)$ forms a regular sequence). We conclude that $X$ is Cohen-Macaulay at all closed points $p$. Since $\dim X = 1$, the depth of $X$ is also 1, so $X$ contains no embedded points, hence reduced. Finally, on a reduced scheme, a function that vanishes everywhere must be the zero function. This completes the proof. 
\end{solution}

\newpage

\section*{FOAG 26.3.G}

\begin{solution}
    Consider the two copies of $\AA^3$ glued together at the origin. This is given by the scheme 
    \[
        X = \Spec R = \Spec \frac{k[x,y,z,a,b,c]}{(ax, bx, cx, ay, by, cy, az, bz, cz)}.
    \]
    We first claim that $X$ is not normal. It suffices to show that at the origin $\pm = (0, 0, 0, 0, 0, 0)$, the localization is not an integral domain. The function $w \in R_{\pm}$ is a zero-divisor, since $aw = 0$. It follows that $R_{\pm}$ is not a domain, hence not integrally closed. 

    Next, we claim that $X$ satisfies $R^1$. If we consider the open subscheme $X - \{0\}$, we simply get two disjoint copies of $\AA^3 - \{0\}$, which is clearly both regular. Also, since the origin has codimension 3, we see that $X$ satisfies $R^1$. Finally, we claim that $X$ satisfies $S^2$ at all points of codimension at most 2. Similarly we may consider $X - \{0\}$, which is clearly normal, so it must satisfy $S^2$ at all points. Since we have removed a codimension-3 point, we see that $X$ satisfies $S^2$ at all points of codimension at most 2. 
\end{solution}

\end{document}